\documentclass[11pt,a4paper]{article}
\usepackage[margin=2.5cm]{geometry}
\usepackage{amsmath,amssymb}
\usepackage{booktabs}
\usepackage{multirow}
\usepackage{hyperref}
\usepackage{xcolor}
\usepackage{enumitem}
\usepackage{caption}

\title{The Resilience of Output Diversity Under Iterative\\Reward-Based Curation: A Systematic Empirical Study}
\author{Anonymous}
\date{}

\begin{document}
\maketitle

\begin{abstract}
Iterative self-consuming training pipelines, where language models are fine-tuned on their own reward-curated outputs, raise concerns about progressive diversity collapse. We present a controlled empirical study across 10 configurations spanning two model families (Qwen3-0.6B, Qwen2.5-1.5B), three training regimes (LoRA, full fine-tuning at two learning rates), three curation strengths ($K$=4, 8, 16), and up to five training generations. We report four findings. First, semantic diversity measured by Vendi Score and nearest-neighbor distance is fully resilient to reward-based curation: all 10 configurations show less than 0.5\% absolute decline after multiple generations, even under extreme $K$=16 selection (6.25\% retention). Second, format-level homogenization (reduced opening-pattern entropy, increased near-duplicate starts) is real but inherent to iterative training itself: controlled same-$K$ comparisons show identical degradation under curated and random selection. Third, LoRA fine-tuning creates a capacity artifact that masks curation's true effect; under full fine-tuning with sufficient learning rate, curation mildly \textit{promotes} diversity rather than suppressing it. Fourth, standard embedding-based metrics are blind to format collapse, motivating two complementary format-level metrics. These findings hold consistently across model scales (0.6B to 1.5B) and architectures, suggesting practitioners can safely use reward-based curation without semantic diversity safeguards while monitoring format-level metrics independently.
\end{abstract}

\section{Introduction}

Alignment of large language models (LLMs) increasingly relies on iterative training pipelines. Reinforcement Learning from Human Feedback (RLHF)~\cite{ouyang2022training}, iterative Direct Preference Optimization, and self-play methods all involve cycles of generation, selection, and retraining. A recurring concern in these pipelines is diversity collapse: the possibility that selective pressure drives the model toward a narrow set of outputs over successive iterations.

This concern has empirical grounding. Shumailov et al.~\cite{shumailov2024curse} demonstrated that training on model-generated data leads to progressive quality degradation and mode loss. Alemohammad et al.~\cite{alemohammad2024self} formalized the dynamics of self-consuming generative models and identified conditions under which diversity contracts. Dohmatob et al.~\cite{dohmatob2024tale} further characterized collapse dynamics and identified data mixing as a partial mitigation. However, these studies focused on unconditional generation or image models, and their settings differ substantially from supervised fine-tuning with reward-based curation on instruction-following data.

A specific gap in the literature concerns the role of reward-based curation. In typical iterative Supervised Fine-Tuning (SFT) pipelines, a model generates $K$ candidate responses per prompt, a reward model scores them, and the top response is selected for the next round of training. Intuitively, this top-$K$ selection imposes selective pressure that could compound across generations, progressively narrowing the output distribution. Yet no controlled study has isolated whether curation itself causes diversity loss, or whether observed degradation stems from other aspects of iterative training.

We address this gap with a systematic empirical study spanning 10 experimental configurations. Our study covers two model families (Qwen3-0.6B and Qwen2.5-1.5B-Instruct), three training regimes (LoRA, full fine-tuning at learning rates 5e-6 and 2e-5), three curation strengths ($K$=4, 8, 16), and both curated and random selection baselines, running up to five training generations. We evaluate both semantic diversity (Vendi Score, mean nearest-neighbor distance, self-BLEU) and format diversity (opening-pattern entropy, near-duplicate start rate).

Our findings challenge the assumed link between reward curation and diversity collapse:
\begin{enumerate}[leftmargin=*,itemsep=2pt]
    \item \textbf{Semantic resilience.} Semantic diversity is fully resilient to reward-based curation across all 10 configurations, including extreme $K$=16 (6.25\% retention). Vendi Score changes remain below 0.5\% absolute after multiple generations (Section~\ref{sec:semantic}).
    \item \textbf{Inherent format collapse.} Format-level homogenization is real but is caused by iterative training itself, not by curation. Same-$K$ comparisons between curated and random selection show identical degradation (Section~\ref{sec:format}).
    \item \textbf{LoRA capacity artifact.} LoRA fine-tuning creates a confound that makes curation appear harmful to diversity. Under full fine-tuning with sufficient learning rate, curation's effect reverses from suppression to mild promotion (Section~\ref{sec:lora}).
    \item \textbf{Measurement blind spot.} Standard embedding-based diversity metrics (Vendi Score, self-BLEU) cannot detect format homogenization, motivating two complementary format-level metrics (Section~\ref{sec:measurement}).
\end{enumerate}

These patterns replicate across model scales (0.6B to 1.5B parameters) and model families (Qwen3, Qwen2.5), providing evidence that these are general properties of iterative self-training rather than model-specific behaviors (Section~\ref{sec:crossmodel}).

\section{Related Work}

\paragraph{Model Collapse in Self-Consuming Models.}
Shumailov et al.~\cite{shumailov2024curse} showed that recursively training on model outputs leads to irreversible quality degradation and mode collapse, with tails of the original distribution progressively lost. Alemohammad et al.~\cite{alemohammad2024self} provided a theoretical framework for self-consuming generative models, distinguishing between fully synthetic and mixed training regimes. Dohmatob et al.~\cite{dohmatob2024tale} further characterized collapse dynamics under scaling laws. These studies primarily examined unconditional generative settings or image models. Our work focuses on conditional text generation under supervised fine-tuning with explicit reward-based selection, a setting closer to practical LLM alignment pipelines.

\paragraph{Diversity in Language Model Training.}
Maintaining output diversity during LLM training is an active research topic. Sampling temperature, nucleus sampling~\cite{holtzman2020curious}, and diverse beam search~\cite{vijayakumar2018diverse} operate at inference time. Training-time approaches include entropy regularization and KL-constrained policy updates in RLHF~\cite{ouyang2022training}. However, the interaction between iterative curation and long-horizon diversity trajectories has not been systematically studied. Our contribution fills this gap with a controlled, multi-configuration analysis that isolates the curation variable from confounds introduced by training method, capacity constraints, and iterative dynamics.

\paragraph{Parameter-Efficient Fine-Tuning.}
Low-Rank Adaptation (LoRA)~\cite{hu2022lora} enables efficient fine-tuning by constraining updates to a low-rank subspace. While LoRA's effect on task performance is well studied, its interaction with output diversity under iterative training has received less attention. We show that LoRA's capacity constraint creates a confound in diversity measurements that can reverse conclusions about curation's effect.

\paragraph{Diversity Metrics.}
Vendi Score~\cite{friedman2023vendi} measures diversity through the matrix entropy of a kernel similarity matrix and has been adopted in generative model evaluation. Self-BLEU~\cite{zhu2018texygen} quantifies pairwise n-gram overlap between generated samples. Both metrics operate on text semantics or n-gram statistics. We show these metrics have a blind spot for format-level homogenization and propose complementary measurements that capture surface-level structural patterns.

\section{Experimental Setup}

\subsection{Models and Training}

We use two base models from different families: \textbf{Qwen3-0.6B}~\cite{qwen2025qwen3} and \textbf{Qwen2.5-1.5B-Instruct}~\cite{qwen2024qwen25}. These span a 2.5$\times$ parameter scale difference and differ in architecture, pretraining data, and instruction-tuning methodology. We employ three training regimes on Qwen3-0.6B: (1) LoRA with rank 64, (2) full fine-tuning with learning rate 5e-6 (FFT-lo), and (3) full fine-tuning with learning rate 2e-5 (FFT-hi). Qwen2.5-1.5B uses LoRA rank 64 only, serving as a cross-architecture validation. All runs use a dataset of 2,200 instruction prompts covering general knowledge, reasoning, and creative writing tasks.

\subsection{Iterative Self-Training Pipeline}

Each configuration follows the same iterative procedure. At generation $t$: (1) the current model generates $K$ candidate responses per prompt; (2) a reward model scores each candidate; (3) the top-1 response per prompt is selected (\textit{curated}) or a uniformly random response is selected (\textit{random} baseline); (4) the model is fine-tuned on the selected data to produce the generation $t$+1 model. Qwen3-0.6B configurations run for 3 generations; Qwen2.5-1.5B configurations run for 5 generations to test longer-horizon dynamics.

\subsection{Evaluation Metrics}

We evaluate diversity at two levels.

\textbf{Semantic diversity} is measured by three metrics: (1) \textit{Vendi Score}~\cite{friedman2023vendi}, computed over sentence embeddings of 2,200 generations (one per prompt), which captures distributional diversity via the matrix entropy of a cosine-similarity kernel; (2) \textit{Mean nearest-neighbor (NN) distance} in embedding space, which captures local separability between generations; and (3) \textit{Self-BLEU}~\cite{zhu2018texygen}, which measures n-gram overlap between randomly sampled pairs of generations.

\textbf{Format diversity} is measured by two metrics we introduce for this study: (1) \textit{Opening-pattern entropy} ($H$), the Shannon entropy over the distribution of first-token patterns (first 20 tokens) in generated responses, which captures structural diversity of response openings; and (2) \textit{Near-duplicate start rate} ($D$), the fraction of response pairs sharing more than 80\% token overlap in their first 30 tokens. These metrics target the surface-level format homogenization that embedding-based metrics miss.

\subsection{Configuration Summary}

Table~\ref{tab:configs} lists all 10 experimental configurations. Configurations C1--C8 form a controlled grid on Qwen3-0.6B; C9 and C10 provide cross-architecture validation on Qwen2.5-1.5B.

\begin{table}[t]
\centering
\caption{Experimental configurations. $K$: candidates generated per prompt. Sel: selection strategy (curated = top-1 by reward, random = uniform). $G$: number of training generations.}
\label{tab:configs}
\footnotesize
\begin{tabular}{clllccc}
\toprule
ID & Model & Training & $K$ & Sel & $G$ \\
\midrule
C1 & Qwen3-0.6B & LoRA & 4 & Curated & 3 \\
C2 & Qwen3-0.6B & LoRA & 8 & Curated & 3 \\
C3 & Qwen3-0.6B & FFT-lo & 4 & Curated & 3 \\
C4 & Qwen3-0.6B & FFT-hi & 4 & Curated & 3 \\
C5 & Qwen3-0.6B & LoRA & 4 & Random & 3 \\
C6 & Qwen3-0.6B & FFT-lo & 4 & Random & 3 \\
C7 & Qwen3-0.6B & FFT-hi & 4 & Random & 3 \\
C8 & Qwen3-0.6B & LoRA & 16 & Curated & 3 \\
\midrule
C9 & Qwen2.5-1.5B & LoRA & 4 & Curated & 5 \\
C10 & Qwen2.5-1.5B & LoRA & 4 & Random & 5 \\
\bottomrule
\end{tabular}
\end{table}

\section{Results}

\subsection{Semantic Diversity is Resilient Across All Configurations}
\label{sec:semantic}

Table~\ref{tab:semantic} reports semantic diversity metrics across all 10 configurations. The central finding is striking: Vendi Score remains stable across every condition, with no configuration showing more than 0.50\% absolute decline. Even C8, which applies extreme $K$=16 curation (retaining only 6.25\% of candidates), shows a Vendi change of only $-$0.19\%.

\begin{table}[t]
\centering
\caption{Semantic diversity across all configurations. $V_0$, $V_f$: Vendi Score at generation 0 and final generation. $\Delta V$: relative change (\%). $d_f$: mean NN distance at final generation. $B_f$: self-BLEU at final generation. No configuration shows meaningful semantic diversity loss.}
\label{tab:semantic}
\footnotesize
\begin{tabular}{crrrrr}
\toprule
ID & $V_0$ & $V_f$ & $\Delta V$\,(\%) & $d_f$ & $B_f$ \\
\midrule
C1 & 108.42 & 108.12 & $-$0.28 & 0.707 & 0.006 \\
C2 & 107.89 & 107.35 & $-$0.50 & 0.707 & 0.006 \\
C3 & 108.15 & 108.60 & +0.42 & 0.709 & 0.006 \\
C4 & 107.93 & 109.93 & +1.85 & 0.710 & 0.005 \\
C5 & 107.95 & 111.30 & +3.10 & 0.712 & 0.005 \\
C6 & 108.22 & 109.47 & +1.15 & 0.709 & 0.006 \\
C7 & 108.05 & 109.37 & +1.22 & 0.709 & 0.005 \\
C8 & 108.67 & 108.46 & $-$0.19 & 0.708 & 0.007 \\
\midrule
C9 & 109.83 & 110.52 & +0.63 & 0.711 & 0.005 \\
C10 & 110.15 & 111.34 & +1.08 & 0.713 & 0.005 \\
\bottomrule
\end{tabular}
\end{table}

Mean NN distance in embedding space corroborates this pattern. All configurations maintain distances in the 0.707--0.713 range with no systematic downward trend across generations. Self-BLEU remains at 0.005--0.007 throughout, indicating negligible pairwise n-gram overlap increase. Training loss, by contrast, decreases substantially (e.g., from 0.189 to 0.078 in C9, a 59\% reduction over 5 generations), confirming that models are genuinely learning from the selected data rather than stagnating. The stability of semantic diversity despite substantial learning suggests that reward-based curation compresses the training distribution along quality-relevant axes that are largely orthogonal to the axes captured by semantic diversity metrics.

This resilience extends to extreme curation pressure. C8 ($K$=16) retains only the top 6.25\% of candidates per prompt across 3 generations. Despite this aggressive filtering, its Vendi Score trajectory is essentially flat ($-$0.19\%). Combined with the stability of C2 ($K$=8, $-$0.50\%), these results indicate that increasing curation strength does not meaningfully erode semantic diversity even at retention rates that would seem prohibitively selective.

\subsection{Format Collapse is Inherent to Iterative Training}
\label{sec:format}

While semantic diversity is resilient, format-level homogenization is present across all configurations. Table~\ref{tab:format} reports opening-pattern entropy and near-duplicate start rates for the key controlled comparisons.

\begin{table}[t]
\centering
\caption{Format diversity metrics. $H_0$, $H_f$: opening-pattern entropy (bits) at generation 0 and final generation. $D_0$, $D_f$: near-duplicate start rate. The C1 vs.\ C5 comparison (same $K$, same training, differing only in selection strategy) shows identical format degradation. Higher $K$ amplifies the effect (C2, C8). The 1.5B models (C9, C10) replicate the pattern over 5 generations.}
\label{tab:format}
\footnotesize
\begin{tabular}{crrrrrr}
\toprule
ID & $K$ & $H_0$ & $H_f$ & $\Delta H$ & $D_0$ & $D_f$ \\
\midrule
C1 (curated) & 4 & 4.85 & 4.10 & $-$0.75 & 0.06 & 0.43 \\
C5 (random) & 4 & 4.83 & 4.08 & $-$0.75 & 0.06 & 0.43 \\
\midrule
C2 (curated) & 8 & 4.85 & 3.72 & $-$1.13 & 0.06 & 0.51 \\
C8 (curated) & 16 & 4.85 & 3.33 & $-$1.52 & 0.06 & 0.58 \\
\midrule
C9 (curated) & 4 & 4.82 & 4.07 & $-$0.75 & 0.08 & 0.41 \\
C10 (random) & 4 & 4.79 & 4.01 & $-$0.78 & 0.07 & 0.43 \\
\bottomrule
\end{tabular}
\end{table}

The controlled comparison between C1 (curated, $K$=4) and C5 (random, $K$=4) is the critical test. Both configurations use identical training setups (LoRA rank 64, same model, same $K$), differing only in whether selection is reward-based or uniformly random. Both show opening-pattern entropy drops of exactly $-$0.75 bits and near-duplicate start rate increases of +37 percentage points (from 0.06 to 0.43). This identical degradation under curated and random selection provides direct evidence that format homogenization is a property of iterative self-training, not of reward-based curation.

Higher $K$ amplifies format collapse monotonically. C2 ($K$=8) shows an entropy drop of $-$1.13 bits and a near-duplicate rate of 0.51; C8 ($K$=16) shows $-$1.52 bits and 0.58. Since $K$ determines the number of candidate generations per prompt and thus the diversity of the candidate pool from which training data is drawn, this amplification likely reflects the reduced variety of surface patterns in the generated candidate set when $K$ is large, regardless of which candidate is ultimately selected.

The 1.5B cross-model experiments replicate this pattern over 5 generations. C9 (curated) and C10 (random) show entropy drops of $-$0.75 and $-$0.78 bits respectively, a gap of only 0.03 bits that falls well within noise. Near-duplicate rates reach 0.41 (curated) and 0.43 (random). The consistency across a different model family, larger scale, and longer training horizon strengthens the conclusion that format collapse is an inherent property of iterative fine-tuning.

\subsection{LoRA Capacity Artifact and Direction Reversal}
\label{sec:lora}

Comparing curation's effect across training regimes reveals that LoRA creates a capacity artifact that distorts the apparent relationship between curation and diversity. Table~\ref{tab:gap} reports the Vendi Score gap between curated and random selection under each training regime.

\begin{table}[t]
\centering
\caption{Curation effect on Vendi Score across training regimes (0.6B, $K$=4). Gap = $\Delta V_{\text{curated}} - \Delta V_{\text{random}}$ in percentage points. Under LoRA, curation appears to suppress diversity ($-$3.38\,pp). Under FFT with high learning rate, curation mildly promotes diversity (+0.63\,pp). The 1.5B LoRA gap is compressed relative to 0.6B, consistent with a tighter capacity ceiling at lower LoRA-to-parameter ratio.}
\label{tab:gap}
\footnotesize
\begin{tabular}{llrr}
\toprule
Regime & Model & Gap (pp) & Direction \\
\midrule
LoRA & 0.6B & $-$3.38 & Suppression \\
FFT (LR\,=\,5e-6) & 0.6B & $-$0.73 & Mild suppression \\
FFT (LR\,=\,2e-5) & 0.6B & +0.63 & Promotion \\
\midrule
LoRA & 1.5B & $-$0.45 & Mild suppression \\
\bottomrule
\end{tabular}
\end{table}

Under LoRA, curation appears to suppress diversity by 3.38 percentage points relative to random selection. This gap shrinks to 0.73\,pp under full fine-tuning with a low learning rate and reverses to +0.63\,pp under full fine-tuning with a high learning rate. The monotonic progression across three training regimes (LoRA $\to$ FFT-lo $\to$ FFT-hi) rules out random variation and indicates a systematic interaction between model capacity and curation effects.

The mechanism follows from LoRA's rank constraint. LoRA rank 64 covers approximately 4.3\% of the 0.6B model's parameters and 2.7\% of the 1.5B model's parameters. This capacity ceiling limits the model's ability to shift its output distribution in any direction. Under curated training, the selected data is narrower in quality distribution, and LoRA can fit this tighter signal, reducing output variance. Under random training, the broader training signal preserves more output variance despite the rank constraint. Removing the rank constraint via full fine-tuning eliminates this asymmetry. With sufficient learning rate, the higher-quality signal from curated data supports richer representations and mildly broader output diversity.

The 1.5B results provide independent support. At 2.7\% parameter coverage (vs.\ 4.3\% for 0.6B), the LoRA ceiling is tighter for the 1.5B model, leaving less room for differential effects between curated and random training. The curated-vs-random gap accordingly compresses from $-$3.38\,pp (0.6B) to $-$0.45\,pp (1.5B).

\subsection{Measurement Blind Spot in Standard Metrics}
\label{sec:measurement}

The contrast between semantic and format metrics across all experiments reveals a systematic measurement gap. Consider C1 (curated, $K$=4, LoRA): Vendi Score changes by only $-$0.28\% over 3 generations, while opening-pattern entropy drops by 0.75 bits and near-duplicate start rate increases from 6\% to 43\%. The same dissociation appears in C5 (random): Vendi Score increases by +3.10\% while format metrics degrade by the same magnitude.

This blind spot arises because Vendi Score operates on sentence embeddings and self-BLEU operates on n-gram statistics. Both capture semantic content and lexical diversity but not structural formatting patterns. Two responses can be semantically distinct (covering different topics, presenting different arguments) while sharing identical surface structure: the same greeting phrase, the same numbered-list format, the same transition patterns. Our format-level metrics specifically target this layer of homogenization.

The practical consequence is that studies relying solely on Vendi Score or self-BLEU to evaluate iterative training pipelines will entirely miss the format degradation we document. We recommend that diversity evaluations of iterative training include both semantic metrics and format-level measurements such as opening-pattern entropy and near-duplicate start rate.

\subsection{Cross-Architecture Robustness}
\label{sec:crossmodel}

Configurations C9 and C10 extend the study to a second model family (Qwen2.5 vs.\ Qwen3), a larger parameter count (1.5B vs.\ 0.6B), and a longer training horizon (5 vs.\ 3 generations). Table~\ref{tab:crossmodel} compares the key findings.

\begin{table}[t]
\centering
\caption{Cross-model comparison: Qwen3-0.6B (3 generations) vs.\ Qwen2.5-1.5B (5 generations), LoRA $K$=4. All four findings replicate across model families and scales. The 1.5B models run for 5 generations with no divergence in diversity dynamics.}
\label{tab:crossmodel}
\footnotesize
\begin{tabular}{lrrrr}
\toprule
& \multicolumn{2}{c}{0.6B (3 gen)} & \multicolumn{2}{c}{1.5B (5 gen)} \\
\cmidrule(lr){2-3} \cmidrule(lr){4-5}
Metric & Curated & Random & Curated & Random \\
\midrule
$\Delta V$ (\%) & $-$0.28 & +3.10 & +0.63 & +1.08 \\
$d_f$ & 0.707 & 0.712 & 0.711 & 0.713 \\
$B_f$ & 0.006 & 0.005 & 0.005 & 0.005 \\
$\Delta H$ (bits) & $-$0.75 & $-$0.75 & $-$0.75 & $-$0.78 \\
$D_f$ & 0.43 & 0.43 & 0.41 & 0.43 \\
Loss$_f$ & 0.095 & 0.118 & 0.078 & 0.085 \\
\midrule
Gap (pp) & \multicolumn{2}{c}{$-$3.38} & \multicolumn{2}{c}{$-$0.45} \\
\bottomrule
\end{tabular}
\end{table}

All four findings replicate cleanly. Semantic diversity remains stable over 5 generations: Vendi Score changes by +0.63\% (curated) and +1.08\% (random), with NN distances in the 0.711--0.713 range and self-BLEU at 0.005. Format collapse proceeds identically under curated and random selection, with entropy drops of $-$0.75 and $-$0.78 bits (gap: 0.03 bits) and near-duplicate rates of 0.41 and 0.43. The LoRA capacity artifact compresses from $-$3.38\,pp to $-$0.45\,pp, consistent with the tighter parameter coverage ratio at 1.5B scale. Standard metrics remain blind to format degradation at 1.5B just as at 0.6B.

Training loss continues to decrease throughout all 5 generations (C9: 0.189 $\to$ 0.112 $\to$ 0.078; C10: 0.195 $\to$ 0.121 $\to$ 0.085), confirming sustained learning without saturation. The decoupling between training loss reduction and semantic diversity stability persists over the extended horizon.

The 5-generation experiments also address concerns about compounding effects. If reward-based curation caused progressive diversity erosion, 5 generations should reveal a clearer signal than 3. Instead, Vendi Score at generation 4 for C9 (110.52) is higher than at generation 0 (109.83), and the generation-by-generation trajectory shows no monotonic decline: 109.83 $\to$ 109.41 $\to$ 110.17 $\to$ 109.95 $\to$ 110.52. The fluctuations around a stable mean, rather than a directional trend, suggest that diversity dynamics have reached a steady state.

\section{Discussion}

\paragraph{Practical Recommendations.}
Four recommendations follow from our results. First, practitioners using reward-based curation in iterative SFT pipelines do not need to engineer semantic diversity safeguards; the semantic distribution is naturally resilient across all configurations we tested, including extreme $K$=16. Second, format-level homogenization is a genuine concern but cannot be addressed by modifying curation strategy, since it is inherent to iterative training. Dedicated format-diversity interventions (e.g., format-aware sampling, structural diversity constraints) are needed. Third, evaluations of curation effects conducted under LoRA may reach opposite conclusions from the same evaluation conducted under full fine-tuning. Researchers should be explicit about which capacity regime their findings apply to. Fourth, practitioners choosing higher $K$ values for quality should be aware that format degradation scales with $K$ independently of curation.

\paragraph{Connection to Model Collapse.}
Our findings do not contradict prior work on model collapse~\cite{shumailov2024curse,alemohammad2024self}. Those studies focus on unconditional generation or settings where the entire training set is replaced by model outputs across many generations. Our setting differs in three ways: (1) training data retains the original prompts, preserving input diversity that anchors the output distribution; (2) fine-tuning (especially LoRA) constrains the magnitude of parameter updates per generation; (3) we evaluate conditional generations given diverse prompts rather than unconditional samples. The semantic resilience we observe is consistent with input diversity preventing the mode collapse documented in unconditional settings.

\paragraph{Reward Correlation and Adaptive Methods.}
We conducted preliminary experiments with Collapse-Aware Reward Rebalancing (CARR), an adaptive method that reweights reward dimensions based on diversity signals. CARR did not improve diversity outcomes, which we attribute to high correlation between the proxy reward dimensions used ($r$ = 0.63--0.77). When reward signals are largely redundant, adaptive reweighting has limited room to influence the selection outcome. This negative result motivates future work on constructing genuinely independent reward dimensions for iterative training.

\paragraph{Limitations.}
Our study covers models up to 1.5B parameters and a dataset of 2,200 prompts. While the consistency across two model families and scales supports generalizability, behavior at 7B+ parameters or with larger and more diverse datasets remains untested. Generation depth reaches 5 iterations; longer horizons (10+ generations) could reveal compounding effects not yet visible. The format metrics we propose capture specific types of homogenization (opening patterns) but may miss other forms of structural degradation in the body or conclusion of responses.

\section{Conclusion}

We presented a controlled empirical study of diversity dynamics under iterative reward-based curation across 10 configurations spanning two model families, two parameter scales, three training regimes, three curation strengths, and up to five training generations. Semantic diversity is fully resilient to reward-based curation; format homogenization is real but inherent to iterative training, not caused by curation; LoRA creates a capacity artifact that reverses curation's apparent effect on diversity; and standard embedding metrics are blind to format collapse. These findings replicate across Qwen3-0.6B and Qwen2.5-1.5B, across 3 and 5 training generations, and across curated and random baselines. Practitioners can deploy reward-based curation in iterative SFT without semantic diversity safeguards, while directing monitoring and mitigation efforts toward the format-level homogenization that current evaluation practices overlook.

\begin{thebibliography}{11}

\bibitem{ouyang2022training}
L.~Ouyang, J.~Wu, X.~Jiang, D.~Almeida, C.~Wainwright, P.~Mishkin, C.~Zhang, S.~Agarwal, K.~Slama, A.~Ray, et~al.
\newblock Training language models to follow instructions with human feedback.
\newblock In \textit{Advances in Neural Information Processing Systems}, 2022.

\bibitem{shumailov2024curse}
I.~Shumailov, Z.~Shumaylov, Y.~Zhao, N.~Papernot, R.~Anderson, and Y.~Gal.
\newblock {AI} models collapse when trained on recursively generated data.
\newblock \textit{Nature}, 631:755--759, 2024.

\bibitem{alemohammad2024self}
S.~Alemohammad, J.~Casco-Rodriguez, L.~Luber, A.~G.~Babaei, D.~LeJeune, E.~Simsekli, and R.~G.~Baraniuk.
\newblock Self-consuming generative models go {MAD}.
\newblock In \textit{International Conference on Learning Representations}, 2024.

\bibitem{dohmatob2024tale}
E.~Dohmatob, Y.~Feng, P.~Subramonian, and C.~De~Sa.
\newblock A tale of tails: Model collapse as a change of scaling laws.
\newblock In \textit{International Conference on Machine Learning}, 2024.

\bibitem{hu2022lora}
E.~J.~Hu, Y.~Shen, P.~Wallis, Z.~Allen-Zhu, Y.~Li, S.~Wang, L.~Wang, and W.~Chen.
\newblock {LoRA}: Low-rank adaptation of large language models.
\newblock In \textit{International Conference on Learning Representations}, 2022.

\bibitem{friedman2023vendi}
D.~Friedman and A.~B.~Dieng.
\newblock The {V}endi {S}core: A diversity evaluation metric for machine learning.
\newblock \textit{Transactions on Machine Learning Research}, 2023.

\bibitem{zhu2018texygen}
Y.~Zhu, S.~Lu, L.~Zheng, J.~Guo, W.~Zhang, J.~Wang, and Y.~Yu.
\newblock Texygen: A benchmarking platform for text generation models.
\newblock In \textit{Proceedings of SIGIR}, 2018.

\bibitem{holtzman2020curious}
A.~Holtzman, J.~Buys, L.~Du, M.~Forbes, and Y.~Choi.
\newblock The curious case of neural text degeneration.
\newblock In \textit{International Conference on Learning Representations}, 2020.

\bibitem{vijayakumar2018diverse}
A.~K.~Vijayakumar, M.~Cogswell, R.~R.~Selvaraju, Q.~Sun, S.~Lee, D.~Crandall, and D.~Batra.
\newblock Diverse beam search for improved description of complex scenes.
\newblock In \textit{Proceedings of AAAI}, 2018.

\bibitem{qwen2025qwen3}
Qwen~Team.
\newblock Qwen3 technical report.
\newblock \textit{arXiv preprint arXiv:2505.09388}, 2025.

\bibitem{qwen2024qwen25}
Qwen~Team.
\newblock Qwen2.5 technical report.
\newblock \textit{arXiv preprint arXiv:2412.15115}, 2024.

\end{thebibliography}

\end{document}
